\section{Task Setting}

The paper studies Tangut short-text translation under two coupled constraints. The first is supervision scarcity: the real data in this repository contains only 491 pairs, split into 400 training examples, 41 development examples, and 50 test examples by a single random shuffle with seed 42. The second is output mismatch: unlike ordinary machine-translation benchmarks, the targets are short and often title-like. Many resemble book titles, sutra titles, or compact scholastic expressions, so success depends on precise lexical choices and compact formatting rather than on generic sentence fluency. These same constraints create the two failure modes addressed in the rest of the paper: local models can overfit noisy supervision and over-generate, while generic MT metrics can reward plausible paraphrase rather than exact title recovery.

This means the task is narrow but meaningful. We ask how to recover short historical titles under only 400 real training examples, not how to solve open-domain Tangut translation. To train local models at all, we therefore use synthetic data as structured proxies. For each SFT branch, 50,000 synthetic Tangut-Chinese pairs are produced from monolingual classical Chinese and dictionary knowledge, then combined with the 400 real training examples upsampled by a factor of 10 to form a 54,000-example training set. All learned models use the same instruction-tuned 7B backbone. These synthetic inputs are not claims about faithful Tangut reconstruction; they are judged only by downstream utility on held-out real Tangut titles.

We also audit data hygiene conservatively. There is zero exact source, target, or full-pair overlap across the train/dev/test splits, zero exact overlap between the 50 test references and the final targets of the 50k synthetic multitask corpus, and zero exact overlap between the frontier few-shot exemplars and the test references. This does not prove the absence of all leakage, but it rules out the simplest exact-copy explanation for the reported gains.

Evaluation is the second bottleneck, so we separate primary evidence from supporting diagnostics. chrF++ is a reasonable overlap anchor for this setting because character-based scoring is relatively stable under orthographic variation \cite{popovic2017chrfpp}. Perplexity is computed as $\exp(\ell)$, where $\ell$ is the mean token-level causal-LM loss under Qwen2.5-0.5B; in this title-heavy setting it becomes extremely large even for the gold reference, so we treat it as diagnostic rather than dispositive. ``Contamination'' in the current paper means mixed-script or artifact contamination in the \emph{output} string, such as Latin letters, literal \texttt{[UNK]} tokens, or prompt residue, not train-test data leakage. Learned MT metrics and QE systems can capture richer signals than surface overlap \cite{rei2020comet, kepler2019openkiwi}, but they can also have domain-specific failure modes \cite{zouhar2024cometpitfalls, pombal2025mint}. We therefore do not use off-the-shelf COMET as a headline arbiter: standard checkpoints are designed for modern sentence-level MT rather than Tangut short-title reconstruction. Accordingly, chrF++, exact match, contamination, suffix preservation, and the reference-aware judge are treated as primary evidence, while lexical coverage, perplexity, and the legacy reference-free judge remain diagnostic.
